\chapter{Point-charge crystal-field estimate (menu 11)}
\label{ch:menu11}
\menulabel{11}

\section{Purpose}

A first, model-level estimate of the crystal-field potential the surrounding
point charges create -- computed from the geometry alone, with no
quantum-chemical calculation. Use it to see the sign pattern, the size and the
symmetry of the field your ligands produce, and as a sanity check on a fitted
parameter set or on a charge assignment.

\section{What you need}

\begin{itemize}
  \item an XYZ structure;
  \item a point charge per element, in units of the elementary charge
    (\code{O=-2,H=0.4} style);
  \item for values in cm$^{-1}$: the radial expectation values
    \code{r2,r4,r6} in \AA$^k$. These are ion- and method-specific and are
    deliberately \emph{not} shipped: without them the estimate stays
    geometry-only, and says so.
\end{itemize}

\section{Walkthrough}

\lstinputlisting[style=fbkcode, firstline=54, lastline=92,
  title={The menu-11 example session (an octahedral CeO\textsubscript{6} toy
  structure, geometry-only)}]
  {examples/expected/m11_point_charge.txt}

\section{What you get}

A report (\file{structure.xyz.fbk.md}) with: the model and its $r < R_i$
assumption, the charges used, the lattice sums per $(k,q)$, the parameters in
cm$^{-1}$ when radial moments are given (otherwise the explicit refusal), the
seven f-orbital energies (eigenvalues of the potential matrix) and the overall
splitting, the symmetry note, and the boundary paragraphs -- including the
measured caveat that point-charge estimates are qualitative.

\section{How to read the numbers}

\begin{readbox}
\begin{itemize}
  \item In the example the only non-zero symmetry-relevant terms are
    $(4,0), (4,4), (6,0), (6,4)$ -- exactly the $O_h$ allowed set, with the
    whole second rank vanishing. That is the check to run first: more non-zero
    terms than your point group allows means the geometry or the charge
    assignment breaks the symmetry.
  \item The seven orbital energies split the 4f shell; with unit radial moments
    their spread ($0.0469$ in the example's units) is a shape measure, not yet
    a spectroscopic splitting. Supply $\langle r^k\rangle$ to get cm$^{-1}$.
  \item These are \emph{$l$-shell} parameters acting inside the 4f orbital
    space. They are not the $J$-manifold Stevens parameters of \menu{10}; the
    operator-equivalent projection is not performed here. Compare the two
    modules through the spectra they generate, never parameter by parameter.
\end{itemize}
\end{readbox}

\section{Worked example}

The session above; then add a charge set that breaks the symmetry (make one
oxygen neutral) and watch the $(2,0)$ and $(2,2)$ lattice sums appear -- the
symmetry reading changes immediately, which is exactly what this estimate is
good for. For quantitative parameters use the ab initio route
(\menu{10} and its cited sources), not this model.

\section{Caveats, refusals and related sections}

\begin{refusalbox}
\begin{itemize}
  \item \textbf{No radial moments, no cm$^{-1}$}: the program refuses to
    produce parameters rather than substitute a default $\langle r^k\rangle$.
    The lattice sums are geometry only.
  \item An element in the structure without a charge is refused, with the
    missing symbols named -- nothing is silently charged as zero.
  \item Atoms sitting on the centre are excluded with a note (the model has no
    self-interaction term).
  \item The model is qualitative: measured deviations between point-charge and
    ab initio crystal fields make it unusable for quantitative extraction
    (the citation is printed in the section). Read it as a geometry and
    charge-assignment check.
\end{itemize}
\end{refusalbox}

Related: \menu{2} (coordination geometry), \menu{10} (the ab initio fit) and
Chapter~\ref{ch:criteria} for the multipole-expansion conventions.
